\section*{Resumen ejecutivo}
\addcontentsline{toc}{section}{Resumen ejecutivo}

El presente informe desarrolla el taller de medición de calidad en uso para el caso MediSalud HIS mediante ISO/IEC 25022. El estudio parte de un conjunto original de \textbf{3.000 incidentes} registrados entre el 2 de enero y el 19 de diciembre de 2025 y lo complementa, sin modificarlo, con \textbf{6.000 eventos operativos} y \textbf{2.000 encuestas simuladas}. La simulación local reproduce tareas clínicas y administrativas con fallos controlados; su finalidad es generar evidencia reproducible, no representar un sistema productivo ni aplicar anticipadamente las mejoras propuestas.

El procedimiento integró los doce escenarios del taller en cinco momentos: comprensión del caso y del modelo SQuaRE; clasificación de incidentes; diseño de atributos y métricas; automatización, visualización e interpretación; y formulación del ciclo de mejora. Los escenarios 1 a 3 se desarrollan con base literal en el Documento Padre. Para los escenarios 4 a 12, cuyo contenido completo no fue suministrado, se reconstruyó una propuesta metodológica coherente con el índice del taller y esta condición se declara en la trazabilidad.

Los resultados automatizados producen diez indicadores. Tres se encuentran en estado verde, cuatro en amarillo y tres en rojo. Los hallazgos prioritarios son: \textbf{41 incidentes de privacidad}, \textbf{4,50\% de errores de facturación} frente a una meta inferior al 1\%, y \textbf{96,52\% de recetas correctas} frente a una expectativa de 100\% por su impacto clínico. También requieren atención el percentil 90 del registro HCE (11,71 segundos), el éxito de agendamiento (87,86\%), el CSAT (3,42 sobre 5) y las caídas de teleconsulta (8,68\%).

Se concluye que MediSalud no puede evaluarse únicamente por disponibilidad técnica. La calidad en uso depende del resultado alcanzado por cada perfil, de los recursos consumidos, de la percepción del usuario, del daño evitable y de la capacidad del sistema para operar en distintas sedes, dispositivos y condiciones de red. El plan de mejora se formula como una fase posterior y mantiene intacta la línea base defectuosa para permitir comparaciones válidas.

\textbf{Palabras clave:} calidad en uso, ISO/IEC 25022, SQuaRE, sistema hospitalario, simulación, métricas, MediSalud HIS.

\vfill
\begin{center}
\begin{tikzpicture}
  \foreach \x/\c/\n/\t in {0/mediteal/3/VERDE,3.7/medigold/4/AMARILLO,7.4/medicoral/3/ROJO}{
    \fill[\c!12] (\x,0) rectangle +(3.2,1.4);
    \fill[\c] (\x,0) rectangle +(.12,1.4);
    \node[anchor=west,font=\fontsize{20}{22}\selectfont\bfseries,text=\c] at (\x+.35,.88) {\n};
    \node[anchor=west,font=\scriptsize\bfseries,text=medigrey] at (\x+.35,.35) {\t};
  }
\end{tikzpicture}
\end{center}
\clearpage
